% TOP_PROBLEM: {"id": 2162, "title": "Erdős Problem #531", "category_id": 3, "category": {"id": 3, "name": "graph_theory", "display_name": "Graph Theory", "description": "Problems involving graphs, networks, and their properties.", "slug": "graph-theory", "order_index": 3, "created_at": "2026-07-31T15:26:25.671Z"}, "status": "open"}
% FIRST_POSTED: null
% ENTRY_KIND: statement_only
% Imported from: unsolvedmath_additional_500.tex; UnsolvedMath ID 2162
% !TeX program = lualatex
% Compile twice: lualatex 2162.tex
% Self-contained: no external bibliography, images, or \input files.
% Numeric section labels and visible numbers are original UnsolvedMath IDs.
\documentclass[11pt,a4paper]{article}
\usepackage[margin=24mm,headheight=14pt]{geometry}
\usepackage{fontspec}
\setmainfont{Latin Modern Roman}
\setsansfont{Latin Modern Sans}
\setmonofont{Latin Modern Mono}
\usepackage{amsmath,amssymb,mathtools}
\usepackage{unicode-math}
\setmathfont{Latin Modern Math}
\usepackage{microtype}
\usepackage{enumitem,xurl,needspace}
\usepackage[dvipsnames]{xcolor}
\usepackage{hyperref}
\hypersetup{unicode=true,colorlinks=true,linkcolor=MidnightBlue,urlcolor=MidnightBlue,
 pdftitle={UnsolvedMath Problem 2162},pdfauthor={Source-annotated compilation},bookmarksnumbered=true}
\usepackage{bookmark,tocloft,titlesec,fancyhdr}
\setlength{\cftsecnumwidth}{4.2em}
\cftsetpnumwidth{2.5em}
\cftsetrmarg{4em}
\titleformat{\section}{\Large\bfseries\raggedright}{\thesection}{0.7em}{}
\pagestyle{fancy}\fancyhf{}
\fancyhead[L]{\small\sffamily Additional UnsolvedMath statements}
\fancyhead[R]{\small Original catalogue IDs}
\fancyfoot[C]{\thepage}
\setlength{\parindent}{0pt}
\setlength{\parskip}{5pt plus 1pt}
\setlength{\emergencystretch}{3em}
\setcounter{tocdepth}{1}\setcounter{secnumdepth}{1}
\setlist[itemize]{leftmargin=3em,itemsep=4pt,topsep=4pt,parsep=0pt}
\allowdisplaybreaks
\newcommand{\N}{\mathbb{N}}
\newcommand{\Z}{\mathbb{Z}}
\newcommand{\Q}{\mathbb{Q}}
\newcommand{\R}{\mathbb{R}}
\newcommand{\C}{\mathbb{C}}
\newcommand{\F}{\mathbb{F}}
\newcommand{\A}{\mathbb{A}}
\newcommand{\PP}{\mathbb{P}}
\newcommand{\E}{\mathbb{E}}
\newcommand{\eps}{\varepsilon}
\newcommand{\dd}{\,\mathrm d}
\newcommand{\sref}[2]{\hyperlink{source:#1:#2}{[S#2]}}
\newif\ifshowcatalogue
\showcataloguetrue
\newcommand{\cataloguescope}[1]{\ifshowcatalogue
 \paragraph{Statement recorded in the supplied source.}
 #1\fi}
\begin{document}
% UnsolvedMath ID 2162; source catalogue code EP-531

\setcounter{section}{2161}

\section{Quantitative Thresholds for Monochromatic Subset Sums}\label{2162}

{\small\sffamily UnsolvedMath ID 2162 · Graph Theory}

\subsection{Definitions and mathematical statement}

\paragraph{Shared notation.}Unless a section says otherwise, $\N=\{1,2,\ldots\}$,
$\Z$ is the ring of integers, $\Q,\R,\C$ are the rational, real, and complex
fields, $[N]=\{1,\ldots,\lfloor N\rfloor\}$, and $\log$ is the natural
logarithm. For real functions of a parameter tending to infinity,
$f\ll g$ and $f=O(g)$ mean $|f|\leq Cg$ eventually for some fixed $C$;
$f\gg g$ reverses this comparison; $f=o(g)$ means $f/g\to0$;
$f\sim g$ means $f/g\to1$; and $f\asymp g$ means both order bounds.
A subscript on $O$ or $\ll$ specifies allowed constant dependence.
The expression $n^{o(1)}$ in an upper bound means growth smaller than
$n^\epsilon$ for each fixed $\epsilon>0$. Local definitions govern any
exception, finite-field convention, or use of the same symbol for another
object. An asymptotic claim is not automatically an effective algorithm.



Every nonempty subset sum of the selected $k$ positive integers must lie in the coloured interval and receive the same colour. Sums can have repeated values, but a selected element cannot be used more than once in one subset sum. \sref{2162}{1} \sref{2162}{2}

\paragraph{Statement recorded in the supplied source.}
Let $F(k)$ be the minimal $N$ such that if we two-colour $\{1,\ldots,N\}$ there is a set $A$ of size $k$ such that all subset sums $\sum_{a\in S}a$ (for $\emptyset\neq S\subseteq A$) are monochromatic. Estimate $F(k)$. \sref{2162}{1}

{\small\emph{Transcription note.} Only unambiguous escape/formatting damage and nonmathematical serialization debris were repaired in the displayed source statement.}

\subsection{Short English statement}

How long must a two-coloured interval be to force a prescribed number of integers whose nonempty subset sums all have one colour?

\subsection{Sources}

{\small

\begin{itemize}

\item[\hypertarget{source:2162:1}{S1}] ULAM.AI, \emph{UnsolvedMath}, supplied \texttt{problems.json}, record 2162 (EP-531), statement and background. The embedded literature-review date is 2026-08-17; that is the archive’s review date, not a fresh certification. Dataset landing page: \url{https://huggingface.co/datasets/ulamai/UnsolvedMath}.

\item[\hypertarget{source:2162:2}{S2}] T. F. Bloom, \emph{Erdős Problems}, Problem 531, \url{https://www.erdosproblems.com/531}. Reference imported from the supplied record; the live problem page was not independently checked for this compilation.

\item[\hypertarget{source:2162:3}{S3}] József Balogh, Sean Eberhard, Bhargav Narayanan, Andrew Treglown, and Adam Zsolt Wagner, Folkman's theorem for random subsets, arXiv:1703.02473 (2017). \url{https://arxiv.org/abs/1703.02473}. Bibliographic information imported from the supplied record.

\item[\hypertarget{source:2162:4}{S4}] Balogh, J\'{o}zsef and Eberhard, Sean and Narayanan, Bhargav and Treglown, Andrew and Wagner, Adam Zsolt, An improved lower bound for Folkman's theorem. Bull. Lond. Math. Soc. (2017), 745-747. Bibliographic information imported from the supplied record.

\item[\hypertarget{source:2162:5}{S5}] Erd\H{o}s, Paul and Spencer, Joel, Monochromatic sumsets. J. Combin. Theory Ser. A (1989), 162-163. Bibliographic information imported from the supplied record.

\end{itemize}

}

\label{end:2162}
\end{document}
